\exercisesection{}

\begin{enumerate}
\def\labelenumi{\arabic{enumi})}
\item
  给出一个群 G、一个域 K 以及 G 的一个忠实表示的例子，使得相应的 \(K[G]\)-模并不忠实。
\item
  令 \(\pi\) 为 G 在希尔伯特空间 E 上的酉表示，且 \(x \in E\) 的范数为 1，并满足
\end{enumerate}

\[
\sup _ {g \in \mathrm{G}} \| \pi (g) x - x \| \leqslant 1.
\]

存在非零的 G-不变向量 \(y \in E\)。

\begin{enumerate}
\def\labelenumi{\arabic{enumi})}
\setcounter{enumi}{2}
\tightlist
\item
  令 E 为希尔伯特空间。记 \(\mathbf{U}(\mathbf{E})\) 为 E 的酉自同态群，记 \(\mathbf{GL}(\mathbf{E})\) 为 E 的可逆自同态群。
\end{enumerate}

\begin{enumerate}
\def\labelenumi{\alph{enumi})}
\item
  群 \(\mathbf{U}(\mathrm{E})\) 赋予由 \(\mathcal{L}(\mathrm{E})\) 上逐点收敛拓扑诱导的拓扑后是拓扑群，记为 \(\mathbf{U}(\mathrm{E})_{s}\)；若 E 是无限维的，则 \(\mathbf{GL}(\mathrm{E})\) 赋予由逐点收敛拓扑诱导的拓扑后不是拓扑群。
\item
  \(\mathbf{U}(\mathrm{E})_{s}\) 的拓扑与由紧收敛拓扑诱导的拓扑一致。
\item
  若 E 是可数生成的，则群 \(\mathbf{U}(\mathrm{E})_{s}\) 可度量化且完备。
\item
  令 L 为由 \(u \in \mathcal{L}(\mathrm{E})\) 且满足 \(u^{*} = -u\) 的元素所成的空间；它是实李群 G = U(E) 的李代数，其中 G 赋予由 \(\mathcal{L}(\mathrm{E})\) 的巴拿赫空间拓扑诱导的拓扑（LIE, III, 第 146 页推论 2）。从 L 到 \(\mathbf{U}(\mathrm{E})_s\) 的指数映射是连续的。
\item
  若 E 是无限维的，则 \(\mathbf{U}(\mathrm{E})_{s}\) 不是型 \(\{L\}\) 的实李群。（考虑 \(\mathbf{U}(\mathrm{E})_{s}\) 中由 E 的某个固定标准正交基下的对角酉自同态组成的子群 K，并注意 K 是紧的。）
\item
  指数映射不是从 L 到 \(\mathbf{U}(\mathbf{E})_{s}\) 的局部同胚。（考虑它在某个固定标准正交基下的对角自同态上的限制。）
\end{enumerate}

\begin{enumerate}
\def\labelenumi{\arabic{enumi})}
\setcounter{enumi}{3}
\tightlist
\item
  令 G 为拓扑群。若存在一个线性泛函 \(\lambda\colon\mathcal{C}_{b}(\mathrm{G};\mathbf{R})\to\mathbf{R}\)，使得
\end{enumerate}

\begin{enumerate}
\def\labelenumi{(\roman{enumi})}
\item
  若 \(f \geqslant 0\)，则 \(\lambda(f) \geqslant 0\)。
\item
  有 \(\lambda(1)=1\)。
\item
  对每个 \(g \in \mathbf{G}\) 及每个 \(f \in \mathcal{C}_b(\mathbf{G};\mathbf{R})\)，有 \(\lambda(\boldsymbol{\gamma}(g)f) = \lambda(\boldsymbol{\delta}(g)f)\)，其中
\end{enumerate}

\[
\boldsymbol {\gamma} (g) f (x) = f \left(g ^ {- 1} x\right), \quad \boldsymbol {\delta} (g) f (x) = f (x g).
\]

若 G 是离散群，此定义与 EVT, IV, 第 73 页习题 4 中的定义一致。

\begin{enumerate}
\def\labelenumi{\alph{enumi})}
\tightlist
\item
  拓扑群 G 可和，当且仅当对于每个整数 \(n \in N\)、每个有限族 \((f_i)_{1 \leqslant i \leqslant n}\)（属于 \(\mathcal{C}_b(\mathbf{G}; \mathbf{R})\)）、G 中的两个族 \((g_i)_{1 \leqslant i \leqslant n}\) 和 \((h_i)_{1 \leqslant i \leqslant n}\) 以及每个实数 a，条件
\end{enumerate}

\[
\sum_ {i = 1} ^ {n} (f _ {i} (x) - f _ {i} (g _ {i} x h _ {i})) \geqslant a
\]

对所有 \(x \in G\) 成立都蕴含 \(a \leqslant 0\)。（为证明该条件充分，考虑形如

\[
x \mapsto \sum_ {i = 1} ^ {n} (f _ {i} (x) - f _ {i} (g _ {i} x h _ {i}))
\]

的函数所成的集合，并应用 Hahn--Banach 定理。）

\begin{enumerate}
\def\labelenumi{\alph{enumi})}
\setcounter{enumi}{1}
\item
  每个紧拓扑群都是可和的，每个阿贝尔拓扑群也都是可和的。
\item
  令 G 为非交换自由群。离散群 G 不是可和的。
\item
  对每个希尔伯特空间 E 及 G 在 E 中的每个连续有界表示，E 上都存在一个希尔伯特范数 q，使得 G 在赋予 q 的希尔伯特空间 E 中的表示是酉的。
\end{enumerate}

\begin{enumerate}
\def\labelenumi{\arabic{enumi})}
\setcounter{enumi}{4}
\tightlist
\item
  令 G 为拓扑群，H 为 G 的开子群。记 p 为典范投影 \(G \rightarrow G/H\)。
\end{enumerate}

\begin{enumerate}
\def\labelenumi{\alph{enumi})}
\item
  令 \(\sigma\colon\mathrm{G/H}\to\mathrm{G}\) 为满足 \(p\circ\sigma=\mathrm{Id}_{\mathrm{G/H}}\) 的映射。记 \(\beta\colon\mathrm{G}\times\mathrm{G/H}\to\mathrm{H}\) 为由 \(\beta(g,x)=\sigma(gx)^{-1}g\) 定义的映射。映射 \(\beta\) 是连续的。
\item
  令 \(\pi\) 为 H 在希尔伯特空间 E 中的有界表示。记 F 为希尔伯特空间 \(\ell^2 (\mathrm{G / H};\mathrm{E})\)。公式
\end{enumerate}

\[
\varrho (g) f (x) = \pi (\beta (g ^ {- 1}, x) ^ {- 1}) f (g ^ {- 1} x)
\]

定义了 G 在 F 中的连续有界表示。

\begin{enumerate}
\def\labelenumi{\alph{enumi})}
\setcounter{enumi}{2}
\tightlist
\item
  若 F 上存在一个希尔伯特范数 q，使得表示 \(\varrho\) 是 G 在 \((F,q)\) 中的酉表示，则 E 上存在一个希尔伯特范数 \(\widetilde{q}\)，使得 \(\pi\) 是 H 在 \((E,\widetilde{q})\) 中的酉表示。
\end{enumerate}

\begin{enumerate}
\def\labelenumi{\arabic{enumi})}
\setcounter{enumi}{5}
\tightlist
\item
  \begin{enumerate}
  \def\labelenumii{\alph{enumii})}
  \tightlist
  \item
    令 E 为复巴拿赫空间，令 \(u \in \mathcal{L}(\mathrm{E})\)。若存在 \(n \in N\)，使得 \(n \geqslant 1\) 且 \(u^{n} = 1_{E}\)，则 u 的谱包含于 n 次单位根群，并由 u 的特征值组成。b) 令 G 为交换拓扑群，且 G 的每个元素都有有限阶。G 在巴拿赫空间中的每个连续不可约表示维数均为 1，且取值于 \(U \subset C^{*}\)——模为 1 的复数所成的群。
  \end{enumerate}
\end{enumerate}

\begin{enumerate}
\def\labelenumi{\alph{enumi})}
\setcounter{enumi}{2}
\tightlist
\item
  令 G 为交换拓扑群，且 G 的有限阶元素集在 G 中稠密。G 在巴拿赫空间中的每个连续不可约表示维数均为 1，且取值于 \(U \subset C^{*}\)——模为 1 的复数所成的群。
\end{enumerate}

\begin{enumerate}
\def\labelenumi{\arabic{enumi})}
\setcounter{enumi}{6}
\tightlist
\item
  对每个整数 \(n \geqslant 1\)，记 \(s_{n}\) 为由下式定义的函数 \(Z \rightarrow R\)
\end{enumerate}

\[
s _ {n} (k) = \left\{ \begin{array}{l l} (1 - k) ^ {n (1 - k)} & \text{若 } k \leqslant - 1 \\ 1 & \text{若 } k = 0 \\ (k + 1) ^ {- (k + 1) / n} & \text{若 } k \geqslant 1. \end{array} \right.
\]

记 S 为函数 \(s_n\) 所成的集合。

\begin{enumerate}
\def\labelenumi{\alph{enumi})}
\item
  对每个整数 \(n \geqslant 1\) 及每个 \((k_{1}, k_{2}) \in \mathbf{Z}^{2}\)，有 \(s_{4n}(k_{1})s_{4n}(k_{2}) \geqslant s_{n}(k_{1} + k_{2})\)。
\item
  对于 \(f \in \mathbf{C}(\mathrm{T})\)，记 \((a_k(f))_{k \in \mathbf{Z}}\) 为 \(f\) 在 0 处的 Laurent 展开的系数。对于每个 \(s \in \mathrm{S}\) 和每个 \(f \in \mathbf{C}(\mathrm{T})\)，级数
\end{enumerate}

\[
\sum_ {k \in \mathbf {Z}} | a _ {k} (f) | s (k)
\]

收敛，并且集合

\[
\mathcal {V} (s, \varepsilon) = \{f \in \mathbf {C} (\mathrm{T}) | \sum_ {k \in \mathbf {Z}} | a _ {k} (f) | s (k) <   \varepsilon \}
\]

对 \(s \in S\) 和 \(\varepsilon > 0\) 定义的这些集合，构成 \(\mathbf{C}(T)\) 上某个局部凸拓扑的 0 的开邻域基本系。乘法 \(\mathbf{C}(T) \times \mathbf{C}(T) \to \mathbf{C}(T)\) 关于此拓扑是连续的。

\begin{enumerate}
\def\labelenumi{\alph{enumi})}
\setcounter{enumi}{2}
\tightlist
\item
  上一问题中定义的拓扑可度量化但不完备。
\end{enumerate}

\(d)\) 对于 \(f \in \mathbf{C}(\mathrm{T})^*\)，记 \(\varrho(f)\) 为映射 \(g \mapsto fg\)，从 \(\mathbf{C}(\mathrm{T})\) 到 \(\mathbf{C}(\mathrm{T})\)。映射 \(\varrho\) 是 \(\mathbf{C}(\mathrm{T})^*\) 在局部凸空间 \(\mathbf{C}(\mathrm{T})\) 中的连续表示。

\begin{enumerate}
\def\labelenumi{\alph{enumi})}
\setcounter{enumi}{4}
\tightlist
\item
  表示 \(\varrho\) 是不可约的。
\end{enumerate}

